\section{Practical Considerations}

%=====================================================================================
\subsection{Ticketmaster}
\label{sec:ticketmaster}
%===========================================================

Key to the good performance of system is an effective way for allocating sequence slots to validators exclusively.
Otherwise, sender broadcasts might ``collide'' in their attempted sequence slots.

It is important to note that in case of collisions, broadcasts are not lost:
\primitive eventually delivers every block \primitivesmall-bcast by a correct sender, albeit without a sequence number. Therefore, eventually a consensus decision will reference it and it will become committed to the total order.  
However, this scenario will suffer the same commit delay as previous protocols, which rely on a leader to commit a wave of accumulated broadcasts; we want to avoid this in \system.
Additionally, eventually blocks carrying sequence numbers need to be enabled for the system to make progress.

To this end, we introduce an abstract functionality referred to as a \emph{ticketmaster}.
The goal of a ticketmaster is to assign sequence slots to actively brodcasting validators, hereon referred to as active validators, exclusively and consecutively.
In order to prevent Byzantine validators from capriciously attempting to \primitivesmall-bcast sequence numbers,
potentially causing collisions and slowdowns,
validators should be allowed to invoke \bbcacast\ only with their sequence slots. 
Validators accept and echo other validators' blocks only when they carry allowed sequence numbers.

There are a number of strategies for realizing a ticketmaster functionality.

\paragraph{Round robin.}
A communication-less strategy is to assign sequence numbers by a round robin rotation among active validators.
For example, say that there are four validators, $p_0$, $p_1$, $p_2$, $p_3$.
By round robin rotation, they will be assigned sequence slots, $0, 1, 2, 3$ respectively, then $4, 5, 6, 7$, and so on.

The advantage of round robin is that it removes contention completely, and that it incurs no extra messages.
The downside is that it requires all validators to be active in each rotation. 
Additionally, any validator that fails to use their slot slows down 
those that come after it because they have to wait for the fallback protocol to finalize the slot.

As a mitigation measure,
the set of active validators may be automatically adjusted every time a fallback consensus decision is made in order to finalize a timed-out slot.
The consensus decision pertaining a new configuration of active validators would have to specify the slot from which the rotation starts.
A new rotation might incur temporary contention due to the switchover, but this does not hurt consistency; it only (potentially) causes a temporary hickup in throughput.

For example, say slot $3$ and $7$ fail to finalize within a designated timeout period.
A fallback consensus can finalize the slots with a $\bot$ value,
and at the same time, will kick validator $p_3$ from the active validators' set
starting at slot number $8$.

To bring $p_3$ back to the system, there would be another fallback consensus decision with a special control command for adding $p_3$ back.


\paragraph{Ticketmaster process.}
Another option is using a ticketmaster server to hand out slot reservations. This works as follows.
When a validator wishes to propose a block, it first reaches out to the ticketmaster process with a ticket-request to be assigned a slot. The ticketmaste responds with a certified slot, which the validator can use to invoke \bbcacast with.
Note that the slot must be certified by the server to prevent Byzantine validators from capriciousy attempting to \bbcacast\ sequence numbers.
Due to the chance of Byzantine faults,
the ticketmaster process should be rotated regularly among nodes every fixed number of slots.

The advantage of this approach is there only active validators need to broadcast.
The downside is that the interaction with a ticketmaster process adds two network latencies to the critical path of \primitive.
However, note that the messages involved in this interaction are very small.
Additionally, batching can amortize the communication cost and delay across a bulk of transmissions.

Experience gained in the benign fault setting~\cite{balakrishnan2012corfu, balakrishnan2020virtual}
indicates that this strategy can perform extremely well,
easily reaching hundreds of thousands to a million commits per second.
The ticketmaster process, albeit being a bottleneck, is very lightweight and stateless, and therefore can perform
with very high throughput.

\paragraph{Hybrid.}

In a hybrid approach, slots are allocated by a rotational regime.
However, each rotation (or group of rotations) is determined by a ticketmaster process on demand.
This works as follows.

Except for the first rotation, each validator sends a ticket request to the ticketmaster process of the next rotation.
The ticketmaster process waits for a certain threshold of requests (a threshold may be determined by assistant policy),
and then \primitivesmall-bcasts a message allocating validators to the next set of rotations.

An advantage of the threshold participation requirement
is guaranteeing a certain degree of participation fairness, namely,
that some fraction of the slots in each rotation will be assigned to correct validators.
This is similar to the full rotation regime, however, the hybrid regime
has the benefit of adaptively determining a rotation by demand.
The downside is the same as having a ticketmaster per request, namely
having a server in the critical path of broadcast deliveries.

\subsection{Transaction dissemination}
Similarly to Narwhal~\cite{Narwhal}, we decouple transaction dissemination from establishing a total order.
Transactions can be disseminated by a consistent broadcast protocol run by worker processes.
Each validator may correspond to multiple worker processes.
Transaction deduplication techniques, as in \cite{Mir} or \cite{RCC} can be applied among workers to ensure effective bandwidth utilisation.
We do not go in depth with respect to transaction dissemination, as this is still a work in progress.

Also, a practical implementation of BBCA would avoid repeatedly sending the actual user payload in \echo, \ready, and \checkpoint messages, but instead with block digests whenever possible.
One can pull the full payload from the validator that first sends a block and a correct validator only needs to do it once throughout the entire \primitive instance. Therefore, in our implementation,
an auxiliary $\langle\textsc{pull}\rangle$ message is used to trigger the resend of an \echo/\ready/\checkpoint message with the full payload in place of the hash. An adversary cannot effectively abuse this pull-style functionality as
a validator needs to resend a full payload at most once to each of other correct validators. It does not impose additional latency when every validators just receive the payload in the \initiate message and also handles the scenario
where the asynchrony of the network could make an \echo/\ready message arrive earlier than \initiate message, with the minimum communication overhead.  Another optimization would be to precautiously send an \abort message upon receiving an \initiate message with a conflicting sequence number.
Last, we follow PBFT~\cite{PBFT} which requires signed messages. However, we could replace signatures with authenticated point-to-point channels, similarly to~\cite{PBFT:2002}.

\subsection{Pacing dissemination via the DAG}
\primitive broadcasts form a DAG, whose primary purpose is the fallback consensus mechanism. 
However,
We can utilize the DAG itself to prevent the use of the DAG for fallback. What we mean by this is the following.

Optimistically, \system should operate at its full speed
even without ever reverting to the  fallback consensus. In theory, one can create contentionless and gap-less
broadcasts.
We are therefore designing \system to synchronize the pace of 
block generations among validators. This is needed because validators would otherwise
keep working on their own slots, where the fastest stays way ahead of others,
leaving lots of gaps which don't contribute to the commit progress. While a
process-based ticketmaster could balance the pace, it is more natural to just
make use of the DAG itself for synchronization. It also helps unify the happy path with
the fallback path as there is already some causality among the blocks. Our current implementation
uses a simple strategy: each validator can at most have $k$ outstanding \primitive instances.
Whenever some instances are finalized, the validator can immediately dispatch new broadcasts. A validator
keeps track of the frontier set (blocks without causal successors), which is used as the references
for a generated block. The frontier is then set to only contain the new block
once it is finalized. This design does not require extra timeouts and
relies on the causality of the DAG to help synchronize the pace.


